\section{Introduction}
\label{sec:intro}

Materials discovery, molecular design and program synthesis increasingly depend on
generative models, and in each the useful output is not merely a plausible one but one that meets a
number: a lattice at a specified density, a composition inside a charge budget, a schedule
within a cost. Learning the distribution of examples does not deliver this, because a model
fit to a corpus reproduces the corpus's spread of the quantity rather than any value a user
picks. Control of a generative model's quantitative behaviour is therefore a prerequisite
for using one as a design tool at all.

The field's standard answer works at the level of whole sequences. Among all distributions
whose expectation of a statistic meets the target, take the one closest in
Kullback--Leibler divergence to the model; the solution is an exponential tilt of the base,
and it is the fixed point that distributional policy gradient and its
descendants pursue. Two properties of that answer are
usually left implicit. It is defined over complete sequences, while an autoregressive model
emits one token at a time, so reaching it requires either sampling from an unnormalized
target or training against an estimator of it.

The second property is the one this paper turns on. An expectation constrains an average,
which says nothing about any individual draw. Constrained optimization has distinguished these
since chance-constrained programming separated expected from per-instance feasibility, and
constrained reinforcement learning inherited the expectation form wholesale. For a generative model the distinction
is not academic, because the samples a user inspects are individual draws rather than
averages, yet the size of the gap between the two notions of feasibility is left
unquantified in this setting.

Working per step is the obvious alternative and has been explored as decoding-time
machinery: grammar and lexical constraints applied as masks, tractable surrogates
consulted during search, and sequential Monte Carlo reweighting partial
sequences. These act on a trained model's output rather than
shaping what the model learns, so a model steered this way was never trained to expect the
steering, and the mismatch between the distribution it learned and the one it is asked to
emit is left for the sampler to absorb.

A separate literature supplies the missing ingredient. Convex programs with implicit
derivatives can sit inside a network as layers, which is what a per-step constraint would
need in order to be trained through rather than bolted on. That
line and the decoding line have not been joined for expectation constraints, so how a
per-step surrogate compares against the sequence-level optimum it approximates has never
been measured against a model actually trained toward that optimum.

Our setting is crystal structure generation, where the target quantity is mass density and
the serialization is text, alongside a synthetic task small enough that every distribution
is computed exactly rather than sampled. Generative models for this domain are established enough that steering one is a fair test
rather than a demonstration; \cref{sec:related} surveys them.

\paragraph{Contributions}

\begin{enumerate}
  \item A differentiable per-step constraint projection with an analytic Jacobian,
        trained through rather than applied at decoding time (\cref{sec:method}).
  \item The concentration result: the surrogate targets a distribution whose
        constrained quantity is concentrated, while the sequence-level optimum
        anti-concentrates it (\cref{sec:concentration}).
  \item Zero-shot retargeting: one checkpoint covers a constraint family with
        statistically bounded equivalence to per-target training
        (\cref{sec:retargeting}).
  \item Four scope boundaries, each measured rather than asserted: a steering
        ceiling, partial and asymmetric two-constraint composition, a serialization
        determinability limit, and inference-time baselines that beat the projection on
        accuracy at a diversity cost (\cref{sec:scope,sec:baselines}).
\end{enumerate}

\begin{kilkey}
The per-step surrogate is not a cheaper approximation of the sequence-level optimum. It
targets a different distribution, one whose constrained quantity is concentrated rather than
spread, and on the tasks measured here that distribution is closer to what a request for a
constraint usually means.
\end{kilkey}

%%%%%%%%%%%%%%%%%%%%%%%%%%%%%%%%%%%%%%%%%%%%%%%%%%%%%%%%%%%%%%%%%%%%%%%%%%%%%%%%
